% SEGMENT OLP-0653-S01
% Part: first-order-logic
% Chapter: model-theory
% Section: nonstandard-arithmetic

\documentclass[../../../include/open-logic-section]{subfiles}

\begin{document}

\olfileid{mod}{bas}{nsa}
\section{எண்கணிதத்தின் நிலையற்ற மாதிரிகள்}

\begin{defn}
$\Lang{L_N}$ என்பது எண்கணிதத்தின் !!{language} ஆகட்டும். அதில்
!!{constant} $\Obj{0}$, இரு இட !!{predicate} $<$, ஓரிட
!!{function} $\prime$, இரு இட !!{function}s $+$ மற்றும் $\times$
ஆகியவை உள்ளன.
\begin{enumerate}
\item எண்கணிதத்தின் \emph{நிலையான மாதிரி} என்பது $\Lang{L_N}$
  க்கான !!{structure}~$\Struct{N}$. அதன் ஆட்களம்
  $\Nat = \{ 0, 1, 2, \dots\}$; அதில் $\Obj{0}$ என்பது $0$ ஆகவும்,
  $<$ என்பது~$\Nat$ மீதான சிறியது-என்னும் தொடர்பாகவும்,
  $\prime$, $+$, $\times$ ஆகியவை முறையே~$\Nat$ மீதான
  அடுத்தெண், கூட்டல், பெருக்கல் செயல்களாகவும் பொருள்கொள்ளப்படுகின்றன.
\item \emph{மெய் எண்கணிதம்} என்பது $\Theory{N}$ கோட்பாடு.
\end{enumerate}
\end{defn}

% SEGMENT OLP-0653-S02
$\Lang{L_N}$ இல் செயல்படும்போது, $\Obj{0}$ இற்கு அடுத்தெண்
செயலியை~$n$ முறைப் பயன்படுத்திப் பெறும்
$\Obj{0}^{\prime\dots\prime}$ வடிவிலான ஒவ்வொரு சொல்லையும்
$\num{n}$ எனச் சுருக்கி எழுதுகிறோம்.

\begin{defn}
$\Lang{L_N}$ க்கான !!{structure}~$\Struct{M}$ ஆனது
$\Struct{N} \iso \Struct{M}$ எனில், எனில் மட்டுமே,
\emph{நிலையானது}.
\end{defn}

% SEGMENT OLP-0653-S03
\begin{thm}\ollabel{thm:non-std}
மெய் எண்கணிதத்திற்கு நிலையற்ற !!{enumerable}
மாதிரிகள் உள்ளன.
\end{thm}

\begin{proof}
புதிய !!{constant}~$c$ ஒன்றைச் சேர்த்து $\Lang{L_N}$ ஐ விரிவாக்கி,
பின்வரும் கோட்பாட்டைக் கருதுக:
\[
\Theory{N} \cup \Setabs{\num{n} < c}{n \in \Nat}.
\]
இக்கோட்பாட்டின் ஒவ்வொரு முடிவுறு பகுதியும் நிறைவுறத்தக்கது.
ஆகவே இறுக்கத்தன்மையால் இதற்கு மாதிரி~$\Struct{M}$ ஒன்று உண்டு;
கீழ்நோக்கிய L\"owenheim--Skolem தேற்றத்தால் அதை
!!{enumerable} ஆக எடுத்துக்கொள்ளலாம்.
$\Domain{M}$ என்பது~$\Struct{M}$ இன் ஆட்களம் என்க. அப்போது
% SOURCE CORRECTION TA-MOD-001: use the defined L_N, and call all interpretations symbols.
$\Lang{L_N}$ இன் தருக்கமல்லாத குறியீடுகளுக்கான
$\Struct{M}$ இன் பொருள்கோள்கள் பின்வருமாறு இருக்கட்டும்:
$\mathbf{z} = \Assign{\Obj{0}}{M} \in \Domain{M}$,
% SOURCE CORRECTION TA-MOD-002: relation on the domain, not on the structure object.
${\prec} =
\Assign{<}{M} \subseteq \Domain{M}^2$,
$* = \Assign{\prime}{M} \colon
\Domain{M} \to \Domain{M}$, மற்றும்
$\oplus = \Assign{+}{M}, \otimes =
\Assign{\times}{M} \colon \Domain{M}^2 \to \Domain{M}$.
ஒவ்வொரு $x \in \Domain{M}$ க்கும், $*$ ஐ~$x$ இற்குப்
பயன்படுத்திப் பெறும் $\Domain{M}$ இன் உறுப்பை $x^*$ என எழுதுகிறோம்.

இப்போது $\Struct{N}$ க்கும் $\Struct{M}$ க்கும் இடையில்
ஓரினச்சார்பு~$h$ ஒன்று உண்டென்றால், ஏதோவொரு
$n \in \Nat$ க்கு $h(n) = \Assign{c}{M}$ ஆக வேண்டும்.
எனவே $s(x) = n$ ஆகும்~$\Struct{N}$ இன் ஏதேனும் ஒதுக்கீடு~$s$
ஐ எடுத்துக்கொள்க. அப்போது
$\Sat{N}{\eq[\num{n}][x]}[s]$; \olref[iso]{thm:isom} இன்
நிரூபிப்பின்படி
$\Sat{M}{\eq[\num{n}][x]}[h \circ s]$ உம் பொருந்தும்.
ஆகவே $\Assign{c}{M} =
\mathbf{z}^{*\cdots *}$ ($*$ ஐ~$n$ முறைத் தொடர்ச்சியாகப்
பயன்படுத்தியது). ஆனால்
$\Sat{M}{\num{n} < c}$ என்றும் $\prec$ தற்சுட்டற்றது
என்றும் கருதப்பட்டதால் இது இயலாது. எனவே $\Struct{M}$
நிலையற்றது.
\end{proof}

% SEGMENT OLP-0653-S04
\begin{prob}
கணம்~$X$ மீதான தொடர்பு~$R$ ஆனது \emph{நன்கு நிறுவப்பட்டது}
என்பது,~$R$ இல் முடிவிலா இறங்கும் சங்கிலிகள் இல்லாதபோது,
அப்போது மட்டுமே; அதாவது
$\dots x_2Rx_1Rx_0$ ஆக அமையும்
$x_0$, $x_1$, $x_2$,~\dots ஆகியவை~$X$ இல் இல்லாதபோது.
செர்மெலோ--ஃபிரேங்கல் கணக் கோட்பாடு~$ZF$ முரணற்றது
என்று கொண்டு, $ZF$ இன் நன்கு நிறுவப்படாத மாதிரிகள்
உள்ளன என்று காட்டுக; அதாவது,
$\dots x_2 \in x_1 \in x_0$ ஆக அமையும்
மாதிரிகள்~$\mathfrak{M}$ உள்ளன என்று காட்டுக.
\end{prob}

% SEGMENT OLP-0653-S05
\olref{thm:non-std} இல் கிடைத்த நிலையற்ற மாதிரி~$\Struct{M}$
நிலையான மாதிரியுடன் அடிப்படைச் சமானமானது. எனவே
$\Struct{M}$ இன் பல பண்புகளை வருவிக்கலாம்.
இந்தப் பிரிவின் எஞ்சிய பகுதி அந்தப் பணியை மேற்கொள்கிறது;
அதன் மூலம் $\Theory{N}$ இன் !!{enumerable} நிலையற்ற
மாதிரிகளுக்குத் துல்லியமான தன்மைப்படுத்தலைப் பெறுவோம்.

\begin{enumerate}
\item $\Domain{M}$ இன் எந்த உறுப்பும் தன்னைவிட
  $\prec$-சிறியது அல்ல: $\lforall[x][\lnot x < x]$
  என்ற வாக்கியம் $\Struct{N}$ இல் மெய்; ஆகவே
  $\Struct{M}$ இலும் மெய்.
\item அதே போன்ற காரணத்தால், $\prec$ ஆனது
  $\Domain{M}$ இன் \emph{நேரியல் வரிசை}; அதாவது
  $\Domain{M}$ மீதான முழுமையான, தற்சுட்டற்ற,
  கடப்புப் பண்புள்ள தொடர்பு.
\item $\mathbf{z}$ என்பது $\Domain{M}$ இன்
  $\prec$-மீச்சிறு உறுப்பு.
\item $\Domain{M}$ இன் ஒவ்வோர் உறுப்பும் அதன்
  $*$-அடுத்தெண்ணைவிட $\prec$-சிறியது; மேலும்
  $x$ ஐவிடப் பெரிய $\Domain{M}$ உறுப்புகளில்
  $\prec$-மீச்சிறியது $x^*$.

% SEGMENT OLP-0653-S06
\item $\Struct{M}$ இல்~$\Nat$ க்கு ஓரினமான
  $\prec$-தொடக்கப் பகுதி ஒன்று உள்ளது:
  $\mathbf{z}, \mathbf{z}^*, \mathbf{z}^{**}, \dots$.
  இதனை $\Domain{M}$ இன் \emph{நிலையான பகுதி}
  என்கிறோம். $\Domain{M}$ இன் மற்ற எந்த உறுப்பும்
  \emph{நிலையற்றது}. நிலையற்ற உறுப்புகள்
  $\Domain{M}$ இல் இருக்கவேண்டும்; இல்லையெனில்
  \olref{thm:non-std} இன் நிரூபிப்பிலுள்ள சார்பு~$h$
  ஓரினச்சார்பாகிவிடும். $\Struct{M}$ இன் இந்நிலையான
  பகுதியில் மதிப்பெடுக்கும் !!{variable}s ஆக
  $n, m, \dots$ ஐப் பயன்படுத்துகிறோம்.
\item ஒவ்வொரு நிலையற்ற உறுப்பும் எந்த நிலையான
  உறுப்பையும்விடப் பெரியது. ஏனெனில்
  ஒவ்வொரு $n \in \Nat$ க்கும்,
  \[
  \Sat{N}{\lforall[z][(\lnot(\eq[z][\Obj{0}] \lor
  \dots \lor \eq[z][\num{n}]) \lif \num{n} < z)]},
  \]
  எனவே $z \in \Domain{M}$ நிலையான உறுப்புகள்
  அனைத்திலிருந்தும் வேறுபட்டால், அவை அனைத்தையும்விட
  அது \emph{பெரியது}.
\item $\mathbf{z}$ தவிர $\Domain{M}$ இன் ஒவ்வோர்
  உறுப்பும் $\Domain{M}$ இன் ஒரேயொரு உறுப்பின்
  $*$-அடுத்தெண்; அந்த முன்னெண்ணை $^*x$
  எனக் குறிக்கிறோம். $x = y^*$ எனில், $x$, $y$
  ஆகியவற்றுள் ஒன்று நிலையானதாக இருந்தால் இரண்டும்
  நிலையானவை; ஒன்று நிலையற்றதாக இருந்தால் இரண்டும்
  நிலையற்றவை.

% SEGMENT OLP-0653-S07
\item $\Domain{M}$ மீதான சமானத் தொடர்பு~$\approx$
  ஐப் பின்வருமாறு வரையறுக்கிறோம்: ஏதோவொரு
  \emph{நிலையான}~$n$ க்கு
  $x \oplus n = y$ அல்லது $y \oplus n =x$
  எனில், எனில் மட்டுமே, $x \approx y$.
  வேறுவிதமாகச் சொன்னால், $x$, $y$ இரண்டும்
  முடிவுறு தொலைவில் இருந்தால், எனில் மட்டுமே,
  $x \approx y$. $n$, $m$ நிலையானவை எனில்
  $n \approx m$. $x$ இன் \emph{தொகுதி}
  என்பது அதன் சமான வகுப்பு
  $[x] = \Setabs{y}{x
  \approx y}$.

% SEGMENT OLP-0653-S08
\item $x \prec y$ மற்றும் $x \not\approx y$
  என்க. $\Sat{N}{\lforall[x][\lforall[y][(x < y \lif (x' < y \lor x' =
      y))]]}$ என்பதால், $x^* \prec y$ அல்லது
  $x^* = y$. இரண்டாவது வாய்ப்பு
  $x \approx y$ என்பதைக் கொடுப்பதால் இயலாது;
  % SOURCE CORRECTION TA-MOD-003: retain the strengthened successor inequality.
  ஆகவே $x^* \prec y$. அதேபோல்,
  $x \prec y$ மற்றும் $x \not\approx y$
  எனில், $x \prec {^*y}$. எனவே
  $x \prec y$ மற்றும் $x \not\approx y$
  எனில், $w \approx x$ ஆகும் ஒவ்வோர்
  உறுப்பும், $v \approx y$ ஆகும் ஒவ்வோர்
  உறுப்பைவிடவும் $\prec$-சிறியது.
  % SOURCE CORRECTION TA-MOD-007: only non-standard blocks extend in both directions.
  இதனால் ஒவ்வொரு நிலையற்ற தொகுதி~$[x]$ உம்
  இரு திசைகளிலும் முடிவிலாத சங்கிலியாக அமைகிறது:
  \[
  \cdots \prec  {^{**}x} \prec {^*}x \prec x \prec x^* \prec x^{**}
  \prec \cdots
  \]
  முழுக்களின் வரிசை வகையைக் கொண்டதால் இது
  $Z$-சங்கிலி எனப்படும்.

% SEGMENT OLP-0653-S09
% SOURCE CORRECTION TA-MOD-004: order only distinct blocks.
\item $\prec$ வரிசையை வெவ்வேறு தொகுதிகளுக்கும்
  உயர்த்தலாம்: $x \prec y$ ஆகவும் அவை
  வெவ்வேறு தொகுதிகளைச் சேர்ந்தவையாகவும் இருந்தால்,
  $x$ இன் தொகுதி $y$ இன் தொகுதியைவிடச் சிறியது.
  நிலையற்ற உறுப்பைக் கொண்ட தொகுதி
  \emph{நிலையற்ற} தொகுதி. நிலையான தொகுதியே
  மீச்சிறு தொகுதி.

% SEGMENT OLP-0653-S10
\item மீச்சிறு நிலையற்ற தொகுதி இல்லை:
  $y$ நிலையற்றது எனில், $x \prec y$,
  $x$ நிலையற்றது, $x \not\approx y$
  என அமையும் உறுப்பு ஒன்று உண்டு.
  நிரூபிப்பு: நிலையான மாதிரி~$\Struct{N}$ இல்
  ஒவ்வோர் எண்ணும் இரண்டால் வகுபடும்; மீதி ஒன்று
  இருக்கலாம். அதாவது
  % SOURCE CORRECTION TA-MOD-005: every y has some half, not every x is its half.
  $\Sat{N}{\lforall[y][\lexists[x][(y = x + x \lor y = x + x +
        \Obj{0}')]]}$.
  அடிப்படைச் சமானத்தால், ஒவ்வொரு
  $y \in \Domain{M}$ க்கும்
  $x \oplus x = y$ அல்லது
  $x \oplus x \oplus \mathbf{z}^*= y$
  ஆக அமையும் $x \in \Domain{M}$ ஒன்று உண்டு.
  $x$ நிலையானதாக இருந்தால் $y$ உம்
  நிலையானதாகிவிடும்; ஆகவே $x$ நிலையற்றது.
  நிலையற்ற உறுப்பு பூச்சியைவிடப் பெரியது
  என்பதால், அது அதன் இரட்டிப்பைவிடச் சிறியது.
  மேலும் $x$, $y$ வெவ்வேறு தொகுதிகளைச்
  சேர்ந்தவை; அதாவது $x \not\approx y$.
  அப்படியல்லாமல், ஏதோவொரு நிலையான~$n$
  க்கு $x \oplus n = y$ என்க.
  $y = x \oplus x$ எனில்,
  $x \oplus n = x \oplus x$;
  கூட்டலின் நீக்கல் விதியால்
  ($\Struct{N}$ இலும், ஆகவே $\Struct{M}$
  இலும் அது பொருந்தும்) $x = n$ கிடைக்கும்.
  அப்போது $x$ நிலையானதாகிவிடும்.
  $y = x \oplus x \oplus \mathbf{z}^*$
  எனும்போதும் இதே போன்றது.
% SEGMENT OLP-0653-S11
\item இதே போன்ற வாதத்தால், மீப்பெரு தொகுதியும் இல்லை.
% SEGMENT OLP-0653-S12
\item தொகுதிகளின் வரிசை செறிந்தது:
  $[x]$ ஆனது~$[y]$ ஐவிடச் சிறிதாகவும்
  ($x \not\approx y$), அவை வேறுபட்டவையாகவும்
  இருந்தால், அவற்றுக்கு இடையில் இரண்டிலிருந்தும்
  வேறுபட்ட தொகுதி~$[z]$ ஒன்று உண்டு.
  $x \prec y$ என்க. முன்போல,
  $x \oplus y$ இரண்டால் வகுபடும்
  (மீதி இருக்கலாம்). எனவே
  $x \oplus y = u \oplus u$ அல்லது
  $x \oplus y = u \oplus u \oplus \mathbf{z}^*$
  ஆக அமையும் $u
  \in \Domain{M}$ ஒன்று உண்டு.
  $u$ என்பது $x$, $y$ ஆகியவற்றின்
  சராசரி; ஆகவே அவற்றுக்கிடையில் உள்ளது.
  $x \oplus y = u \oplus u$ எனக் கொள்க
  (மற்ற நேர்வும் இதே போன்றது).
  $u \approx x$ எனில், ஏதோவொரு
  நிலையான~$n$ க்கு
  \[
  x \oplus y = x \oplus n \oplus x \oplus n,
  \]
  ஆகவே $y = x \oplus n \oplus n$;
  இது கருதுகோளுக்கு மாறாக $x \approx y$
  என்பதைக் கொடுக்கும். எனவே
  $u \not\approx x$. இதே போன்ற வாதம்
  $u \not\approx y$ என்பதையும் தருகிறது.
\end{enumerate}
% SEGMENT OLP-0653-S13
எனவே நிலையற்ற தொகுதிகள் விகிதமுறு எண்களைப்
போல வரிசைப்படுத்தப்பட்டுள்ளன: அவை முனைகளற்ற
!!a{enumerable} செறிந்த நேரியல் வரிசையாக அமைகின்றன.
இதனால் மெய் எண்கணிதத்தின் எந்த இரு
!!{enumerable} நிலையற்ற மாதிரிகள்
$\Struct{M}_1$, $\Struct{M_2}$ ஆகியவற்றின்
$<$ மற்றும் $=$ ஐ மட்டுமே கொண்ட மொழிக்கான
குறைப்புகள் ஓரினமானவை என்று பின்வருகிறது.
உண்மையில் ஓரினச்சார்பு~$h$ ஐப்
பின்வருமாறு வரையறுக்கலாம்:
$\Struct{M_1}$, $\Struct{M_2}$
இன் நிலையான பகுதிகள் நிலையான மாதிரி
$\Struct{N}$ க்கு ஓரினமானவை;
எனவே ஒன்றுக்கொன்றும் ஓரினமானவை.
நிலையற்ற பகுதியை உருவாக்கும் தொகுதிகள்
விகிதமுறு எண்களைப் போலவே வரிசைப்படுத்தப்பட்டுள்ளன;
எனவே \olref[dlo]{thm:cantorQ} படி
அவையும் ஓரினமானவை. தொகுதிகளுக்கிடையிலான
ஓரினச்சார்பை, ஒவ்வொரு தொகுதியிலும்
தன்னிச்சையான உறுப்புகளைப் பொருத்தி,
பின்னர் $\Struct{M_1}$ இல்~$x$
இன் அடுத்தெண்ணின் படிமம்,
$\Struct{M_2}$ இல்~$x$ இன் படிமத்தின்
அடுத்தெண்ணாக இருக்குமாறு அமைத்து,
தொகுதிகளுக்கு \emph{உள்ளேயும்}
ஓரினச்சார்பாக நீட்டிக்கலாம்.
இதனால் $\mathfrak{M}_1$,
$\mathfrak{M}_2$ ஆகியவை எண்கணிதத்தின்
முழு மொழியில் ஓரினமானவை என்று
\emph{பின்வராது} என்பதைக் கவனிக்க:
ஓரினத்தன்மை எப்போதும் ஒரு
குறியீட்டமைப்பைச் சார்ந்தது;
$\Domain{M_1}$, $\Domain{M_2}$
மீதான கூட்டலையும் பெருக்கலையும்
ஓரினமற்ற வழிகளில் வரையறுக்கலாம்.
(ஒரு நிறைவான கோட்பாட்டின்
எண்ணத்தக்க மாதிரிகளின் எண்ணிக்கை~2
ஆக இருக்க முடியாது என்ற வாட்டின்
புகழ்பெற்ற தேற்றத்திலிருந்தும்
இது பின்வருகிறது.)

% SEGMENT OLP-0653-S14
\begin{prob}
எண்கணிதத்தின் நிலையற்ற மாதிரியில்
மீப்பெரு தொகுதி இருக்க முடியாது என்று காட்டுக.
\end{prob}

% SEGMENT OLP-0653-S15
\begin{prob}
$\Lang{L}$ என்பது $\eq$ தவிர
$<$ ஐ மட்டுமே !!{predicate}
ஆகக் கொண்ட முதல்தர !!{language}
ஆகட்டும்; $\Struct{N} = (\Nat,
<)$ என்க. $\Struct{N}$ இன்
முடிவுறும் அல்லது நிரப்பி முடிவுறும்
உட்கணங்கள் அனைத்தும் வரையறுக்கத்தக்கவை.
$\Struct{N}$ இன் வரையறுக்கத்தக்க
உட்கணங்கள் இவை \emph{மட்டுமே}
என்று காட்டுக.

(குறிப்பு: முதலில்
$\Obj{prc}(x,y)$ என்பது
“$x$ ஆனது $y$ இன் உடனடி முன்னெண்”
என்பதைச் சுருக்கி எழுதும்
$\Lang{L}$-வாய்பாடு என்க:
\[
x<y \land \lnot \lexists[z][(x<z \land z < y)].
\]
இப்போது $\Struct{N}$ இன் ஒவ்வொரு
வரையறுக்கத்தக்க உட்கணத்துக்கும்
$\Lang{L}$ இல் வாய்பாடு~$!A(x)$
ஒன்று உண்டு. அத்தகைய ஒவ்வொரு
$!A$ க்கும் வாக்கியம்~$!D$ ஐக் கருதுக:
\[
\lexists[x][\lforall[y][\lforall[z][((x<y \land x<z \land \Obj{prc}(y,z)
  \land !A(y)) \lif !A(z))]]].
\]
$!A$ வரையறுக்கும் $\Struct{N}$
இன் உட்கணம் முடிவுறுவதோ அல்லது
அதன் நிரப்பி முடிவுறுவதோ,
$\Sat{N}{!D}$ எனும்போது,
அப்போது மட்டுமே என்று காட்டுக.

இப்போது $\Struct{N}$ உடன்
அடிப்படைச் சமானமான நிலையற்ற
மாதிரி~$\Struct{M}$ ஐ எடுத்துக்கொள்க.
$a \in \Domain{M}$ நிலையற்றது எனில்,
$a$ ஐவிடப் பெரிய
$b, c \in \Domain{M}$ ஐ எடுத்துக்கொள்க;
$c$ இன் உடனடி முன்னெண்~$b$
ஆகட்டும். அப்போது $h(b)=c$
% SOURCE CORRECTION TA-MOD-006: h is an automorphism of the order reduct.
ஆக அமையும் $\Domain{M}$ இன்
வரிசைச் சுயஓரினச்சார்பு~$h$
ஒன்று உண்டு (ஏன்?).
எனவே $b$ ஆனது~$!A$ ஐ
நிறைவுறுத்தினால் $c$ உம்
நிறைவுறுத்தும் (ஏன்?).
இதனால் $!D$ ஆனது
$\Struct{M}$ இல் மெய்;
ஆகவே $\Struct{N}$ இலும் மெய்.
இதன்படி $!A$ வரையறுக்கும்
$\Struct{N}$ இன் உட்கணம்
முடிவுறுவதோ அல்லது
அதன் நிரப்பி முடிவுறுவதோ
ஆகும்.
\end{prob}

\end{document}
